\documentclass[12pt,a4paper]{article}
\usepackage[utf8]{inputenc}
\usepackage[T1]{fontenc}
\usepackage{geometry}
\geometry{margin=1in}
\usepackage{graphicx}
\usepackage{hyperref}
\usepackage{enumitem}
\usepackage{tabularx}
\usepackage{fancyhdr}
\usepackage{listings}
\usepackage{xcolor}

\lstset{
    basicstyle=\ttfamily\footnotesize,
    breaklines=true,
    captionpos=b,
    numbers=left,
    numberstyle=\tiny,
    frame=single,
    backgroundcolor=\color{gray!10}
}

\overfullrule=0pt

\title{AISAC Documentation}
\author{AISAC Team}
\date{\today}

\begin{document}

\maketitle
\tableofcontents
\newpage

\section{Overview}

\subsection{Welcome to AISAC}

AISAC (AI-led Sensors and Control) is a comprehensive security alert management system designed to receive, process, and manage security alerts from security panels using the SIA DC-09 protocol. The system provides real-time monitoring, intelligent alert routing, and powerful analytics for security operations.

\subsubsection{Key Capabilities}
\begin{itemize}
\item Real-time SIA DC-09 protocol message processing
\item TCP/UDP server for receiving security panel alerts
\item Multi-site and multi-area support
\item Role-based access control (Guards, Heads, Admins)
\item Interactive floor plans and site mapping
\item Intelligent alert assignment and escalation
\end{itemize}

\subsubsection{System Architecture}
\begin{itemize}
\item React + Vite + TypeScript
\item Convex (Real-time database)
\item Node.js TCP/UDP server
\item SIA DC-09 (DC-09-1998.10)
\item Tailwind CSS + shadcn/ui
\end{itemize}

\section{Getting Started}

\subsection{System Requirements}

\subsubsection{Prerequisites}
\begin{itemize}
\item Node.js 18.x or higher (for backend and Convex)
\item Modern web browser (Chrome, Firefox, Edge, Safari)
\item Windows server (x64)
\item Open ports for SIA receiver (default: 7800) and web server (default: 3000)
\end{itemize}

\subsubsection{Hardware Requirements}
\begin{itemize}
\item \textbf{CPU:} Dual-core processor or better
\item \textbf{RAM:} 4GB minimum (8GB recommended)
\item \textbf{Storage:} 2GB available disk space
\item \textbf{Network:} Stable LAN connection
\end{itemize}

\subsection{On-Premises Installation \& Setup}

How to deploy AISAC for on-premises use.

\subsubsection{Step 1: Extract the Deployment Package}
\begin{itemize}
\item Extract the provided ZIP file containing all necessary files to your server.
\item The package includes the frontend build, backend code, and startup scripts.
\end{itemize}

\subsubsection{Step 2: Start the Backend Services}
Double-click the \texttt{start-backend.bat} file to start the Convex backend and SIA receiver:

\begin{lstlisting}[language=bash]
@echo off
REM Start Convex backend
start /B npx convex dev --no-interactive
REM Wait for Convex to start (adjust timeout as needed)
PING 127.0.0.1 -n 6 >NUL
REM Start SIA receiver
start /B node server\server.js
echo AISAC backend and SIA receiver started.
pause
\end{lstlisting}

This will run in the background. A command prompt window will appear briefly and then close.

\subsubsection{Step 3: Start the Frontend}
Double-click the \texttt{start-frontend.bat} file to start the web server:

\begin{lstlisting}[language=bash]
@echo off
REM Start frontend web server
npx serve dist -s -l 3000
pause
\end{lstlisting}

This will start a local web server on port 3000. Open your browser and navigate to \texttt{http://localhost:3000}

\subsubsection{Step 4: Access the Application}
\begin{itemize}
\item Open your web browser and go to \texttt{http://localhost:3000}
\item Log in using the default credentials provided below.
\end{itemize}

\subsubsection{Step 5: (Optional) Seed Initial Data}
\begin{itemize}
\item If you need demo data, run the provided seed scripts by double-clicking \texttt{seed-data.bat}
\item Default user accounts are listed below.
\end{itemize}

\subsection{Default User Accounts}

Use these credentials to log in after setup.

\textbf{Note:} These are demo user accounts provided for testing and demonstration purposes.

\subsubsection{Admin Account}
\begin{tabular}{|l|l|}
\hline
\textbf{Username} & admin \\
\textbf{Password} & admin123 \\
\hline
\end{tabular}

\subsubsection{Head Account}
\begin{tabular}{|l|l|}
\hline
\textbf{Username} & head1 \\
\textbf{Password} & head123 \\
\hline
\end{tabular}

\subsubsection{Guard Accounts}
\begin{tabular}{|l|l|}
\hline
\textbf{Username} & guard1 \\
\textbf{Password} & guard123 \\
\hline
\textbf{Username} & guard2 \\
\textbf{Password} & guard123 \\
\hline
\end{tabular}

\section{Sensor Setup}

Setting up sensors and floor plans.

\subsection{Understanding the Hierarchy}
\begin{itemize}
\item \textbf{Sites (Accounts):} Top-level customer locations, identified by account number (e.g., "3333")
\item \textbf{Areas (Floors/Partitions):} Sub-divisions within a site, identified by area number (e.g., "01", "02")
\item \textbf{Sensors (Zones/Points):} Individual detection devices, identified by zone number (e.g., "0008")
\end{itemize}

\subsection{Creating a New Account}
\begin{enumerate}
\item \textbf{Navigate to Admin Panel:} Only available to Admin users
\item \textbf{Add Account Information:} Enter account number, name, address, and description
\item \textbf{Configure Areas:} Add areas/floors with area numbers matching your security panel configuration
\item \textbf{Upload Floor Plans:} Upload image files (PNG, JPG) for each area
\item \textbf{Place Sensors:} Click on floor plan to add sensors with correct zone numbers
\end{enumerate}

\subsection{Floor Plan Best Practices}
\begin{itemize}
\item Use high-resolution images (1920x1080 or higher recommended)
\item Ensure floor plans are properly scaled and oriented
\item Label sensor positions accurately to match physical installation
\item Use consistent naming conventions for sensors (e.g., "Main Entrance Door", "Server Room Motion")
\item Test sensor positioning by generating test alerts
\end{itemize}

\subsection{Sensor Configuration}

\subsubsection{Required Fields:}
\begin{itemize}
\item \textbf{Name:} Descriptive name for the sensor
\item \textbf{Type:} door, motion, fire, panic, camera, etc.
\item \textbf{Zone Number:} Must match security panel configuration
\item \textbf{Position:} X,Y coordinates on floor plan
\end{itemize}

\textbf{Important:} Zone numbers must exactly match your security panel configuration. Mismatched zone numbers will result in alerts not being linked to the correct sensors.

\section{Features}

Comprehensive security alert management capabilities.

\subsection{Core Features}

\subsubsection{Real-Time Alert Processing}
\begin{itemize}
\item Receives SIA DC-09 formatted messages via TCP/UDP
\item Automatic event classification and prioritization
\item Support for burglary, fire, panic, and access control events
\item Real-time dashboard updates
\item Alert status tracking (unassigned, assigned, in-progress, resolved)
\end{itemize}

\subsubsection{Role-Based Access Control}
\begin{itemize}
\item \textbf{Guard:} View assigned alerts, update status, respond to incidents
\item \textbf{Head:} Assign alerts, monitor teams, view all sites
\item \textbf{Admin:} Full system access, user management, configuration
\item Multi-account access for cross-site monitoring
\end{itemize}

\subsubsection{Interactive Site Mapping}
\begin{itemize}
\item Upload and annotate floor plans
\item Visual sensor placement on floor plans
\item Real-time alert location visualization
\item Multi-area/floor support per site
\item Click sensors to view details and history
\end{itemize}

\subsubsection{Advanced Analytics}
\begin{itemize}
\item Alert trends and patterns over time
\item Response time analytics
\item Event category distribution
\item Peak hours and hotspot analysis
\item Guard performance metrics
\end{itemize}

\subsection{Alert Workflow Features}

\subsubsection{Intelligent Assignment}
Automatically or manually assign alerts to available guards based on priority, location, and workload.

\subsubsection{Guard Availability Status}
Guards can mark themselves as available/away. System prevents assignment to unavailable guards.

\subsubsection{Response Actions}
Predefined response actions: Lockdown, Dispatch, Investigate with guided workflows.

\subsubsection{Notes and Documentation}
Add detailed notes to each alert for audit trails and handover documentation.

\section{User Guide}

Role-specific guides for different user types.

\subsection{Using AISAC by Role}

\subsubsection{Guard User Guide}
Guards are the front-line responders who handle assigned security alerts.

\paragraph{1. Setting Your Availability}
Use the availability toggle at the top of the dashboard to indicate if you're on duty.

\paragraph{2. Viewing Assigned Alerts}
Your dashboard shows alerts assigned to you.

\paragraph{3. Responding to Alerts}
Follow the guided workflow when responding.

\paragraph{4. Using Floor Plans}
Visual floor plans show sensor locations. Click on sensors to view their status and recent alerts.

\subsubsection{Head User Guide}
Heads supervise guards, assign alerts, and monitor overall security operations.

\paragraph{1. Monitoring All Alerts}
View all alerts across your assigned sites.

\paragraph{2. Assigning Alerts}
Assign alerts to available guards based on priority and location.

\paragraph{3. Monitoring Performance}
Track guard response times and resolution rates.

\paragraph{4. Managing Teams}
View team availability and workload distribution.

\subsubsection{Admin User Guide}
Admins have full system access for configuration, user management, and testing.

\paragraph{1. System Configuration}
Configure SIA receivers, database settings, and system parameters.

\paragraph{2. User Management}
Create, edit, and manage user accounts and permissions.

\paragraph{3. Site and Sensor Setup}
Configure sites, areas, sensors, and floor plans.

\paragraph{4. Testing and Maintenance}
Run system tests, view logs, and perform maintenance tasks.

\section{Alert Management}

Understanding Security Alerts - SIA DC-09 protocol event types and priorities.

\subsection{Alert Priorities}

\begin{tabular}{|l|l|}
\hline
\textbf{Priority} & \textbf{Description} \\
\hline
Critical & Immediate response required. Includes fire alarms, panic buttons, hold-up alarms. \\
\hline
High & Urgent attention needed. Includes burglary alarms, forced entry, unauthorized access. \\
\hline
Medium & Requires attention soon. Includes access denied, late to close, supervision issues. \\
\hline
Low & Informational. Includes system restores, routine opens/closes, test messages. \\
\hline
\end{tabular}

\subsection{Common Event Codes (SIA DC-09)}

\begin{tabular}{|l|l|}
\hline
\textbf{Code} & \textbf{Description} \\
\hline
BA & Burglary \\
\hline
FA & Fire Alarm \\
\hline
PA & Panic Alarm \\
\hline
HA & Hold-up \\
\hline
BR & Burglary Restore \\
\hline
OP & Open \\
\hline
CL & Close \\
\hline
TA & Tamper \\
\hline
\end{tabular}

\subsection{Alert Lifecycle}

\begin{enumerate}
\item \textbf{Alert Received:} SIA DC-09 message received via TCP/UDP, parsed, and stored in database with timestamp and priority.
\item \textbf{Assignment:} Head/Admin assigns alert to available guard. Status changes to "assigned".
\item \textbf{In Progress:} Guard accepts assignment and begins response. Can choose lockdown, dispatch, or investigate actions.
\item \textbf{Resolved:} Guard marks alert as resolved with notes. Resolution time and actions are recorded for analytics.
\end{enumerate}

\section{Analytics}

Understanding your security metrics.

\subsection{Analytics Dashboard}

\begin{itemize}
\item \textbf{Alert Trends:} View alert volume over time periods (daily, weekly, monthly). Identify patterns and unusual spikes.
\item \textbf{Response Time Analysis:} Track average time from alert receipt to acknowledgment, assignment, and resolution.
\item \textbf{Event Category Distribution:} See breakdown of alert types: burglary, fire, access control, etc. Identify which systems trigger most alerts.
\item \textbf{Peak Hours Analysis:} Identify times of day with highest alert volume. Optimize guard scheduling accordingly.
\item \textbf{Hotspot Zones:} Identify sensors/zones with most frequent alerts. May indicate equipment issues or high-risk areas.
\item \textbf{Guard Performance:} Track individual guard metrics: alerts handled, average response time, resolution rate.
\end{itemize}

\subsection{Key Performance Indicators}

\begin{itemize}
\item \textbf{Alert Resolution Rate:} Percentage of alerts resolved vs. total alerts received. Target: >95\%
\item \textbf{Average Response Time:} Time from alert receipt to guard action. Target: Critical <5min, High <15min
\item \textbf{First-Time Resolution Rate:} Alerts resolved without escalation or reassignment. Target: >85\%
\item \textbf{Guard Availability:} Percentage of time guards are marked as available. Monitor for adequate coverage.
\end{itemize}

\section{Administration}

Advanced configuration and management.

\subsection{SIA Receiver Configuration}

The SIA receiver runs as a separate Node.js service that can operate in two modes: Server mode (listens for connections) or Client mode (connects to a remote service).

\subsubsection{Server Mode (Default)}
Listens for incoming connections from security panels:

\begin{lstlisting}[language=bash]
# Start in server mode:
npm run server

# Default Configuration:
TCP_HOST: 0.0.0.0
TCP_PORT: 7800
Protocol: SIA DC-09
\end{lstlisting}

Configure your security panels to send messages to the server's IP address on port 7800.

\subsubsection{Client Mode}
Connects to a remote SIA DC-09 service:

\begin{lstlisting}[language=bash]
# Start in client mode (Windows):
$env:CONNECTION_MODE="client"; $env:REMOTE_HOST="127.0.0.1"; $env:REMOTE_PORT="7800"; npm run server

# Start in client mode (Linux/Mac):
CONNECTION_MODE=client REMOTE_HOST=127.0.0.1 REMOTE_PORT=7800 npm run server
\end{lstlisting}

Client mode is useful when connecting to existing SIA receivers or for testing with remote systems.

\subsubsection{Environment Variables}
\begin{itemize}
\item CONNECTION\_MODE: "server" or "client" (default: server)
\item TCP\_HOST: Server mode bind address (default: 0.0.0.0)
\item TCP\_PORT: Server mode listen port (default: 7800)
\item REMOTE\_HOST: Client mode remote IP (default: 127.0.0.1)
\item REMOTE\_PORT: Client mode remote port (default: 7800)
\item CONVEX\_URL: Convex backend URL
\end{itemize}

\subsection{Database Management}

\subsubsection{Convex Database}
The system uses Convex as the real-time database backend.
\begin{itemize}
\item Access Convex dashboard for raw data queries
\item Run migrations using npx convex run migrations
\item Export data for backup purposes
\item Monitor database performance and usage
\end{itemize}

\subsubsection{Database Management CLI}
Interactive CLI tool for managing alerts database:

\begin{lstlisting}[language=bash]
# Start the database manager:
tsx server/dbManager.ts
\end{lstlisting}

Available operations:
\begin{itemize}
\item Count alerts in database
\item List recent alerts (last 20)
\item Migrate old format alerts to SIA DC-09
\item Clear ALL alerts (with confirmation)
\end{itemize}

\subsubsection{Seeding Test Data}
Commands for resetting or seeding database:

\begin{lstlisting}[language=bash]
# Seed initial data:
npx convex run seed

# Seed area map:
npx convex run seedAreaMap

# Clear and reseed:
npx convex run clearAndSeed
\end{lstlisting}

\subsection{Troubleshooting}

\subsubsection{Alerts Not Appearing}
\begin{itemize}
\item Verify SIA receiver is running (npm run server)
\item Check security panel is sending to correct IP/port
\item Review server console logs for parsing errors
\item Confirm account numbers match between panel and database
\end{itemize}

\subsubsection{Sensors Not Linking to Alerts}
\begin{itemize}
\item Verify zone numbers match exactly (remove leading zeros)
\item Check sensor is assigned to correct account number
\item Ensure area numbers match receiver ID format
\end{itemize}

\subsubsection{Guards Not Receiving Assignments}
\begin{itemize}
\item Confirm guard is marked as "Available"
\item Check guard has correct customer account associations
\item Verify user role is set to "guard"
\end{itemize}

\subsubsection{Floor Plans Not Displaying}
\begin{itemize}
\item Verify image was uploaded successfully
\item Check browser console for CORS or loading errors
\item Ensure image format is supported (PNG, JPG)
\end{itemize}

\subsection{Support \& Resources}

\begin{itemize}
\item \textbf{SIA DC-09 Specification:} DC-09-1998.10
\item \textbf{Framework:} React 19 + Vite
\item \textbf{Database:} Convex
\item \textbf{Version:} 1.0.0
\end{itemize}

\end{document}